\chapter{CƠ SỞ LÝ THUYẾT}
\label{chuong2}


Chương này trình bày các nền tảng lý thuyết làm cơ sở cho bài toán ước tính PM2.5 dựa trên vệ tinh được giải quyết trong luận văn này. Mục 2.1 xuất phát từ kỳ vọng ban đầu của tài liệu khoa học --- rằng mô hình đạt R\textsuperscript{2} > 0,8 --- và cho thấy kỳ vọng đó chỉ phản ánh bài toán nội suy thời gian, trong khi bài toán vận hành thực sự là ngoại suy không gian, vốn khó hơn một cách cơ bản. Mục 2.2 tổng quan các kết quả nghiên cứu hiện có và xác định các khoảng trống làm động lực cho công trình này. Mục 2.3 giải thích mối quan hệ vật lý giữa độ sâu quang học sol khí và PM2.5 ở mức mặt đất, chú trọng đến các yếu tố làm phức tạp mối quan hệ này trong môi trường nhiệt đới. Mục 2.4 giới thiệu khung học máy được sử dụng trong luận văn, tập trung vào XGBoost và chiến lược chính quy hóa DART. Mục 2.5 trình bày các phương pháp đánh giá --- đánh giá chéo ngẫu nhiên cho nội suy thời gian, LOSO cho ngoại suy không gian, và đánh giá độc lập như phép kiểm tra chuyển giao.


\section{Từ kỳ vọng ban đầu đến bài toán trọng tâm}

Ước tính PM2.5 mặt đất từ dữ liệu vệ tinh và khí tượng thường được trình bày trong tài liệu khoa học như một bài toán đơn nhất, với các mô hình báo cáo R\textsuperscript{2} > 0,8 như bằng chứng rằng vấn đề đã được giải quyết. Kỳ vọng tự nhiên là một mô hình tương tự, áp dụng tại Việt Nam, sẽ đạt hiệu suất tương đương. Mục này cho thấy kỳ vọng đó dựa trên một nhầm lẫn giữa hai bài toán có bản chất khác nhau.

\subsection{Kỳ vọng ban đầu: nội suy thời gian}

Câu hỏi xuất phát là đơn giản: liệu mô hình kết hợp dữ liệu vệ tinh và khí tượng có thể tái dựng nồng độ PM2.5 mặt đất? Với một trạm có lịch sử quan trắc, mô hình lấp đầy các khoảng trống dữ liệu (cảm biến ngừng hoạt động, bảo trì định kỳ, truyền thông bị gián đoạn) bên trong giai đoạn quan trắc. Mô hình có quyền truy cập chuỗi thời gian của chính trạm đó --- mức trung bình, biên độ ngày đêm, hành vi theo mùa, phản ứng với điều kiện khí tượng --- cả trước và sau khoảng trống. Đây về bản chất là bài toán \emph{nội suy thời gian}, và đánh giá chéo k-fold ngẫu nhiên --- thực hành chủ đạo trong tài liệu khoa học --- là phương pháp đánh giá phù hợp cho nó. Các mô hình thường xuyên đạt R\textsuperscript{2} > 0,7 bởi vì danh tính của trạm cung cấp một tiên nghiệm mạnh: mô hình ngầm biết ``đây là loại trạm nào'' và chỉ cần dự đoán các độ lệch quanh một đường cơ sở đã biết.

Kết quả R\textsuperscript{2} > 0,8 mà tài liệu khoa học báo cáo là hoàn toàn hợp lệ --- cho bài toán nội suy thời gian. Nhưng ý nghĩa vận hành lớn nhất tại Việt Nam không phải lấp khoảng trống tại các trạm đã có, mà là ước tính PM2.5 tại những nơi \emph{không có trạm}. Khi ta đặt câu hỏi đó, bài toán thay đổi hoàn toàn.

\subsection{Bài toán trọng tâm: ngoại suy không gian}

Bài toán thực sự mà luận văn này giải quyết là ngoại suy không gian: ước tính PM2.5 tại một vị trí \emph{chưa từng} được quan trắc. Mô hình phải suy luận không chỉ các dao động theo thời gian mà còn cả mức ô nhiễm nền --- vốn phụ thuộc vào các yếu tố cực kỳ địa phương như mật độ giao thông, mức độ gần các nguồn công nghiệp, sự giam giữ do địa hình, và hình thái đô thị mà vệ tinh chỉ quan sát gián tiếp ở độ phân giải 1--10 km. Mô hình không hề có thông tin lịch sử nào tại vị trí mục tiêu --- nó không biết ``đây là loại vị trí nào'' và phải suy luận điều đó hoàn toàn từ các quan trắc vệ tinh và các trạm lân cận. Đánh giá bỏ-một-trạm-ra (LOSO) là phương pháp đánh giá phù hợp: tất cả dữ liệu của trạm đánh giá bị giữ lại, buộc mô hình hoạt động như thể trạm đó không tồn tại.

Với chỉ 40 trạm quan trắc trên 331.000 km\textsuperscript{2}, mục đích trọng tâm của một mô hình dựa trên vệ tinh tại Việt Nam chính là ước tính PM2.5 ở những vị trí không có trạm --- ở các khu dân cư, trường học, vùng nông thôn nằm cách xa mạng lưới quan trắc. Luận văn này chọn ngoại suy không gian làm bài toán trọng tâm và sử dụng LOSO làm phương pháp đánh giá chính xuyên suốt.

\subsection{Khoảng cách giữa kỳ vọng và thực tế}

Nội suy thời gian và ngoại suy không gian chia sẻ cùng dữ liệu đầu vào và cùng thuật toán, nhưng chúng đánh giá hai khả năng hoàn toàn khác nhau: nội suy thời gian đánh giá khả năng nắm bắt động học theo giờ tại các trạm đã biết, ngoại suy không gian đánh giá khả năng tổng quát hóa đến các vị trí mới. Khi tài liệu khoa học báo cáo một con số R\textsuperscript{2} duy nhất từ đánh giá chéo ngẫu nhiên mà không nêu rõ rằng đó là hiệu suất nội suy thời gian, người đọc dễ hiểu nhầm rằng mô hình có thể dự đoán tốt ở bất kỳ đâu --- bao gồm cả các vị trí mới. Khoảng cách giữa hai con số (Chương 5 cho thấy R\textsuperscript{2} giảm từ khoảng 0,80 xuống khoảng 0,20) chính là nội dung định lượng trọng tâm của luận văn: nó đo chi phí của việc không biết ``đây là loại vị trí nào'' và xác định rằng thông tin về bản sắc trạm --- mức ô nhiễm nền, chế độ theo mùa, vi mô hình đô thị --- là nguồn kỹ năng chính mà ngoại suy không gian phải bù đắp.


\section{Các kết quả nghiên cứu tương tự}

\subsection{Học máy cho ước tính PM2.5 từ dữ liệu vệ tinh}

Việc áp dụng học máy vào ước tính PM2.5 dựa trên vệ tinh đã tăng tốc kể từ khoảng năm 2015, được thúc đẩy bởi sự sẵn có của các sản phẩm vệ tinh độ phân giải cao và các phương pháp tổ hợp ngày càng mạnh mẽ. Hướng tiếp cận chủ đạo bao gồm việc huấn luyện một mô hình có giám sát --- thường là random forest, gradient boosting (XGBoost, LightGBM), hoặc mạng nơ-ron sâu --- để dự đoán PM2.5 đo tại mặt đất, sử dụng kết hợp độ sâu quang học sol khí truy xuất từ vệ tinh, các biến tái phân tích khí tượng (nhiệt độ, độ ẩm, gió, độ cao lớp biên hành tinh), các chỉ báo sử dụng đất, và trong một số trường hợp, đầu ra mô hình vận chuyển hóa học làm ước tính tiên nghiệm.

Tại Trung Quốc, nơi hơn 1.500 trạm quan trắc quốc gia cung cấp dữ liệu huấn luyện dồi dào, Wei et al. đã phát triển mô hình Random Forest Không-Thời gian (STRF) đạt R\textsuperscript{2} > 0,85 ở độ phân giải 1 km. Hướng tiếp cận của họ tận dụng mạng lưới trạm dày đặc để nắm bắt tính bất đồng nhất không gian ở quy mô tinh, và các phương pháp random forest không-thời gian liên quan sau này đã làm nền tảng cho bộ dữ liệu GlobalHighAirPollutants (GHAP), vốn đã trở thành một sản phẩm PM2.5 toàn cầu được sử dụng rộng rãi. Tuy nhiên, như được bàn ở Mục 2.5, R\textsuperscript{2} được báo cáo phản ánh đánh giá chéo dựa trên mẫu, trong đó từng cặp trạm-ngày riêng lẻ được gán ngẫu nhiên vào các fold bất kể danh tính của trạm.

Tại Hoa Kỳ, Di et al. và các nhóm liên quan đã sử dụng các phương pháp tổ hợp mạng nơ-ron với dữ liệu quan trắc EPA, báo cáo R\textsuperscript{2} đánh giá chéo từ 0,76--0,89 tùy theo độ phân giải thời gian. Các nghiên cứu ở châu Âu của de Hoogh et al. đã áp dụng các mô hình lai giữa hồi quy sử dụng đất và học máy với R\textsuperscript{2} đánh giá chéo ở mức trung bình từ 0,53--0,72, mặc dù các nghiên cứu này nhìn chung sử dụng mạng lưới dày hơn của Việt Nam.

Đối với Ấn Độ, Kawano et al. \cite{kawano2025} đại diện cho chuẩn mực vàng về phương pháp luận hiện nay. Sử dụng khoảng 1.000 trạm với XGBoost và các đặc trưng từ MODIS AOD, TROPOMI, ERA5, và GEOS-Chem, họ báo cáo R\textsuperscript{2} = 0,85 dưới đánh giá chéo ngẫu nhiên nhưng R\textsuperscript{2} = 0,67 dưới đánh giá chéo không gian 10-fold (within-R\textsuperscript{2} = 0,49), và tuyên bố rõ ràng rằng đánh giá chéo ngẫu nhiên "phóng đại hiệu suất mô hình trong các ứng dụng thực tế then chốt."

\subsection{Các nghiên cứu PM2.5 tại Việt Nam và Đông Nam Á}

Tài liệu khoa học về mô hình hóa PM2.5 riêng cho Việt Nam còn thưa thớt. Ngo et al. \cite{ngo2023} đã tạo ra một bộ dữ liệu PM2.5 theo ngày cho Việt Nam từ 2012 đến 2020 bằng một mô hình hiệu ứng hỗn hợp với đầu vào MODIS AOD, MERRA-2, và ERA5, được huấn luyện trên 9 trạm mặt đất. Họ báo cáo R\textsuperscript{2} = 0,75 (Pearson r = 0,87), nhưng việc đánh giá sử dụng đánh giá chéo 10-fold dựa trên mẫu. Số lượng trạm ít --- 9 trạm trên toàn quốc --- đồng nghĩa với việc ngay cả các mẫu thời gian của một trạm duy nhất cũng có thể chi phối hiệu suất đánh giá chéo.

Tại Thái Lan, Gupta et al. \cite{gupta2021} đã sử dụng MERRA-2 với random forest trên 51 trạm, báo cáo R = 0,95 dưới đánh giá chéo ngẫu nhiên 10-fold. Các nghiên cứu ở lưu vực sông Mekong của Lu et al. \cite{lu2024} đã áp dụng các mô hình stacking để ước tính PM2.5, nhưng ghi nhận rõ ràng rằng độ chính xác của các sản phẩm PM2.5 toàn cầu "ở lưu vực sông Mekong cần được kiểm chứng thêm."

Một hạn chế chung của các nghiên cứu ở Đông Nam Á là sự phụ thuộc vào đánh giá dựa trên mẫu và sự vắng mặt của đánh giá chéo không gian nghiêm ngặt. Chưa có nghiên cứu nào được công bố thực hiện đánh giá bỏ-một-trạm-ra trên một bộ dữ liệu PM2.5 của Việt Nam, và đây chính là khoảng trống mà luận văn này giải quyết.

\subsection{Các sản phẩm mô hình vận chuyển hóa học}

Các mô hình vận chuyển hóa học (CTM) mô phỏng thành phần khí quyển bằng cách giải các phương trình liên kết về phát thải, vận chuyển, biến đổi hóa học, và lắng đọng trên các miền lưới. Hai sản phẩm PM2.5 dẫn xuất từ CTM thường được sử dụng làm đặc trưng hoặc làm mốc so sánh trong ước tính dựa trên vệ tinh:

GEOS-CF (GEOS Composition Forecasting), được phát triển bởi Văn phòng Mô hình hóa và Đồng hóa Toàn cầu của NASA, cung cấp ước tính PM2.5 toàn cầu theo giờ ở độ phân giải khoảng 25 km. Nó sử dụng cơ chế hóa học GEOS-Chem kết hợp với mô hình khí quyển GEOS. Mặc dù GEOS-CF đã được chứng minh là nắm bắt được các sự kiện ô nhiễm ở quy mô synop, các kiểm kê phát thải của nó cho Đông Nam Á --- đặc biệt với các nguồn di động (xe máy), đốt sinh hoạt (nấu nướng, sưởi ấm), và đốt nông nghiệp (rơm rạ) --- được biết là chưa đầy đủ.

MERRA-2 (Modern-Era Retrospective analysis for Research and Applications, Version 2) cung cấp tái phân tích sol khí có tích hợp các quan trắc AOD vệ tinh thông qua đồng hóa dữ liệu. Một đánh giá toàn cầu gần đây của Gupta et al. \cite{gupta2026} báo cáo Chỉ số Phù hợp chỉ 0,39 cho Đông Nam Á --- hiệu suất kém nhất trong tất cả các khu vực toàn cầu được đánh giá. MERRA-2 đã được ghi nhận là không thể tái tạo biến thiên PM2.5 theo ngày đêm ngay cả ở những khu vực được mô tả đặc trưng kỹ như Đồng bằng Hoa Bắc, với sai số lớn nhất xảy ra trong các đợt ô nhiễm mùa đông.

Bộ dữ liệu GlobalHighAirPollutants (GHAP) cung cấp ước tính PM2.5 toàn cầu theo ngày ở độ phân giải 1 km bằng cách hợp nhất AOD vệ tinh, mô phỏng CTM, và các phép đo mặt đất thông qua một random forest không-thời gian. R\textsuperscript{2} đánh giá chéo nổi bật 0,91 của nó là dựa trên mẫu, và độ chính xác của nó tại Việt Nam chưa được kiểm chứng độc lập. Luận văn này cung cấp đánh giá độc lập đầu tiên như vậy.


\section{Mối quan hệ AOD--PM2.5}

Độ sâu quang học sol khí (AOD) là độ suy giảm tích phân theo cột của bức xạ mặt trời do các sol khí trong khí quyển, được đo dưới dạng một đại lượng không thứ nguyên. Bởi vì PM2.5 là một thành phần của cột sol khí, tồn tại một mối quan hệ vật lý giữa AOD và nồng độ PM2.5 ở mức mặt đất. Tuy nhiên, mối quan hệ này không hề đơn giản, và việc hiểu các giới hạn của nó là thiết yếu để diễn giải các ước tính PM2.5 dựa trên vệ tinh.

Mối quan hệ cơ bản có thể được biểu đạt một cách khái niệm như sau: PM2.5 tỷ lệ với AOD nhân với một hệ số phụ thuộc vào phân bố thẳng đứng của sol khí trong cột khí quyển, tỷ số giữa khối lượng PM2.5 khô và tổng độ suy giảm sol khí (hiệu suất suy giảm theo khối lượng), và hệ số tăng trưởng hút ẩm phản ánh sự hấp thụ nước của các hạt sol khí ở độ ẩm môi trường. Mỗi yếu tố trong số này đưa vào sự biến thiên làm suy yếu tương quan AOD-PM2.5.

Yếu tố phân bố thẳng đứng chủ yếu được chi phối bởi độ cao lớp biên hành tinh (PBLH). Khi lớp biên dày (thường trong điều kiện đối lưu ban ngày), sol khí được phân bố qua một thể tích cột lớn hơn, và một giá trị AOD cho trước tương ứng với nồng độ bề mặt thấp hơn. Ngược lại, trong điều kiện ổn định ban đêm hoặc nghịch nhiệt, sol khí tập trung gần bề mặt, và cùng giá trị AOD đó tương ứng với PM2.5 bề mặt cao hơn. Đây là lý do PBLH nằm trong số những đặc trưng khí tượng quan trọng nhất trong các mô hình dự đoán PM2.5.

Hiệu suất suy giảm theo khối lượng phụ thuộc vào thành phần và phân bố kích thước của sol khí. Các hạt mịn (PM2.5) tán xạ ánh sáng hiệu quả hơn trên mỗi đơn vị khối lượng so với các hạt thô, nhưng thành phần thay đổi theo nguồn: sol khí giàu sulfat (từ nhà máy điện), carbon hữu cơ (từ đốt sinh khối), carbon đen (từ đốt diesel), và bụi khoáng đều có hiệu suất suy giảm khác nhau. Tại Việt Nam, hỗn hợp sol khí bị chi phối bởi khí thải xe máy, đốt nấu nướng và sưởi ấm, các quá trình công nghiệp, và đốt rơm rạ theo mùa, tạo ra một hồ sơ thành phần khác với các điều kiện ở Đông Á và Bắc Mỹ, nơi phần lớn các mô hình AOD-PM2.5 đã được phát triển.

Hệ số tăng trưởng hút ẩm đặc biệt quan trọng ở Việt Nam nhiệt đới, nơi độ ẩm tương đối môi trường thường xuyên vượt 70\% và có thể duy trì trên 85\% trong những giai đoạn kéo dài vào mùa gió mùa và mùa lạnh. Ở độ ẩm cao, các hạt sol khí hút ẩm hấp thụ nước và trương nở, làm tăng tiết diện quang học của chúng và do đó làm tăng AOD mà không có sự gia tăng tương ứng về khối lượng PM2.5 khô. Điều này tạo ra một độ thiên lệch dương theo độ ẩm: các ước tính PM2.5 dựa trên AOD có xu hướng dự đoán vượt trong điều kiện ẩm trừ khi được hiệu chỉnh tường minh. Các cảm biến PM2.5 tán xạ laser chi phí thấp (như Plantower PMS5003 được dùng trong mạng lưới quan trắc của Việt Nam) cũng chịu ảnh hưởng của hiện tượng này, bởi chúng đo tán xạ quang học bao gồm cả đóng góp của lớp vỏ nước quanh hạt.

Bản thân việc truy xuất AOD từ vệ tinh cũng bị phức tạp hóa bởi mây. Thuật toán MODIS MAIAC không thể truy xuất AOD xuyên qua mây, và khí hậu gió mùa của Việt Nam tạo ra độ che phủ mây dai dẳng trong phần lớn thời gian trong năm. Trong bộ dữ liệu của luận văn này, AOD vệ tinh chỉ khả dụng cho khoảng 15\% số quan trắc theo giờ. Dữ liệu thiếu không được phân bố ngẫu nhiên --- điều kiện trời quang có xu hướng trùng với không khí khô hơn, ít ô nhiễm hơn vào mùa hè và với sự đình trệ áp cao (thường ô nhiễm nặng) vào mùa đông --- tạo ra một độ thiên lệch lấy mẫu hệ thống mà mô hình phải thích ứng thông qua việc xử lý giá trị thiếu thay vì nội suy.


\section{XGBoost và chính quy hóa DART}

XGBoost (eXtreme Gradient Boosting) là một thuật toán cây quyết định tăng cường gradient đã trở thành một phương pháp chủ đạo cho các bài toán học có giám sát trên dữ liệu bảng. Nó xây dựng một tổ hợp các cây quyết định một cách tuần tự, với mỗi cây được huấn luyện để dự đoán các sai số phần dư của tổ hợp đứng trước. Dự đoán cuối cùng là tổng đầu ra của tất cả các cây. XGBoost tích hợp một số cơ chế chính quy hóa --- phạt L1 và L2 trên trọng số lá, ràng buộc độ sâu cây tối đa, trọng số con tối thiểu, và lấy mẫu con theo hàng và cột --- giúp giảm quá khớp và cải thiện khả năng tổng quát hóa.

Đối với luận văn này, biến thể DART (Dropouts meet Multiple Additive Regression Trees) của XGBoost được sử dụng. DART giải quyết một điểm yếu cụ thể của gradient boosting tiêu chuẩn: xu hướng các cây ban đầu chi phối tổ hợp trong khi các cây về sau chỉ thực hiện các hiệu chỉnh không đáng kể, một hiện tượng được gọi là chuyên biệt hóa quá mức. Trong DART, ở mỗi vòng lặp boosting, một tập con ngẫu nhiên các cây hiện có bị "loại bỏ" (tạm thời gỡ khỏi tổ hợp), và cây mới được huấn luyện để bù trừ cho đóng góp của các cây bị loại bỏ. Điều này có tác dụng phân bổ tải dự đoán đồng đều hơn trên toàn tổ hợp và giảm tương quan giữa các cây riêng lẻ. Trong bối cảnh dự đoán không gian, nơi mô hình phải tổng quát hóa đến các vị trí chưa thấy với phân bố đặc trưng có thể khác biệt, tổ hợp cân bằng hơn của DART mang lại độ bền tốt hơn so với boosting tiêu chuẩn.

Việc lựa chọn XGBoost thay vì các phương án học sâu (mạng nơ-ron tích chập, transformer, mạng nơ-ron đồ thị) được thúc đẩy bởi ba cân nhắc thực tiễn. Thứ nhất, XGBoost xử lý giá trị thiếu một cách bản địa bằng cách học hướng chia tối ưu cho các đặc trưng thiếu trong quá trình huấn luyện --- điều thiết yếu khi 85\% số quan trắc theo giờ thiếu AOD vệ tinh. Thứ hai, XGBoost cung cấp điểm số tầm quan trọng đặc trưng cho phép diễn giải những gì mô hình đã học, vốn rất quan trọng để hiểu vì sao dự đoán không gian thất bại tại một số trạm nhất định. Thứ ba, hiệu quả tính toán của XGBoost cho phép thực hiện các thí nghiệm loại bỏ và đánh giá chéo rộng khắp (hơn 30 cấu hình, mỗi cấu hình cần 37--40 fold LOSO) vốn cấu thành phần cốt lõi thực nghiệm của luận văn này.


\section{Các phương pháp đánh giá tương ứng}

Mục 2.1 đã phân biệt giữa nội suy thời gian (kỳ vọng ban đầu) và ngoại suy không gian (bài toán trọng tâm). Mỗi bài toán đòi hỏi một phương pháp đánh giá riêng, sao cho phương pháp đánh giá phản ánh đúng bài toán mà người sử dụng quan tâm. Mục này trình bày ba khung đánh giá được sử dụng trong luận văn và liên kết mỗi khung với bài toán mà nó đo lường.

\subsection{Đánh giá chéo k-fold ngẫu nhiên --- đánh giá nội suy thời gian}

Trong đánh giá chéo k-fold ngẫu nhiên, bộ dữ liệu gồm N quan trắc được phân hoạch ngẫu nhiên thành k fold xấp xỉ bằng nhau, không quan tâm đến danh tính trạm hay thứ tự thời gian. Mô hình được huấn luyện trên k-1 fold và được đánh giá trên fold giữ lại, luân phiên qua tất cả k fold. Đối với dữ liệu PM2.5 theo giờ từ S trạm trong T bước thời gian (N = S $\times$ T), việc phân hoạch ngẫu nhiên phân tán các quan trắc của mỗi trạm ra khắp tất cả các fold, nên ở mỗi vòng huấn luyện, mô hình có quyền truy cập khoảng (k-1)/k chuỗi thời gian của mỗi trạm.

Đây chính xác là phương pháp đánh giá phù hợp cho bài toán nội suy thời gian (Mục 2.1.1): nó đo khả năng lấp đầy các khoảng trống dữ liệu tại các trạm có lịch sử quan trắc. Các mô hình thường đạt R\textsuperscript{2} > 0,7 trong chế độ này bởi vì mô hình ngầm học được mức nền và hành vi theo mùa của mỗi trạm từ các giờ huấn luyện. Con số này phản ánh kỹ năng thực sự trong bài toán mà nó đánh giá --- nhưng nó không đo lường khả năng dự đoán tại các vị trí mới.

\subsection{Đánh giá chéo bỏ-một-trạm-ra (LOSO) --- đánh giá ngoại suy không gian}

Trong đánh giá chéo LOSO, tất cả các quan trắc từ một trạm được giữ lại đồng thời, mô hình được huấn luyện trên S-1 trạm còn lại, và các dự đoán được tạo ra cho trạm giữ lại. Quá trình này được lặp lại cho từng trạm trong S trạm, và các tham số hiệu suất được tính cả theo từng trạm lẫn theo tổng hợp.

Đây là phương pháp đánh giá phù hợp cho bài toán ngoại suy không gian (Mục 2.1.2): mô hình phải dự đoán tại một vị trí mà nó chưa từng nhìn thấy bất kỳ quan trắc nào, mô phỏng chính xác tình huống triển khai thực tế. Các đặc trưng nội suy không gian (RFSI) được tính lại trong mỗi fold để loại trừ trạm giữ lại, ngăn chặn rò rỉ không gian trực tiếp.

Hiệu suất LOSO được dự kiến thấp hơn đáng kể so với đánh giá chéo ngẫu nhiên, đặc biệt khi (i) mạng lưới trạm thưa thớt, (ii) các trạm trải dài qua các chế độ ô nhiễm bất đồng nhất, và (iii) trung bình PM2.5 theo trạm biến thiên rộng trên toàn mạng lưới. Cả ba điều kiện đều đúng đối với Việt Nam. Sự chênh lệch giữa hai con số không nên được hiểu là ``đánh giá chéo ngẫu nhiên bị phóng đại'' --- nó phản ánh chính xác rằng hai phương pháp đo lường hai bài toán khác nhau, và bài toán ngoại suy không gian khó hơn bài toán nội suy thời gian một cách cơ bản.

Hai biến thể tổng hợp được sử dụng trong luận văn này. R\textsuperscript{2} trung bình theo trạm có trọng số-n xử lý các dự đoán của mỗi trạm theo tỷ lệ với số quan trắc và phương sai của nó, tóm tắt hiệu suất trên toàn mạng lưới theo cách bị chi phối bởi các trạm có khối lượng lớn, phương sai cao; đây là bản tóm tắt LOSO liên trạm, và nó khác biệt với R\textsuperscript{2} gộp thực sự được tính bằng cách nối phần dư theo giờ của tất cả các trạm thành một vectơ duy nhất (một thống kê dành riêng cho định nghĩa nghiêm ngặt đó và được bàn ở Chương 5). R\textsuperscript{2} trung bình theo trạm tính R\textsuperscript{2} độc lập cho từng trạm rồi lấy trung bình, gán trọng số bằng nhau cho mỗi trạm bất kể khối lượng dữ liệu hay phương sai. Hai tham số có thể phân kỳ đáng kể: trung bình có trọng số-n bị chi phối bởi các trạm phương sai cao (vốn có xu hướng là những trạm ô nhiễm nhất và dễ dự đoán nhất), trong khi R\textsuperscript{2} trung bình theo trạm cho thấy liệu mô hình có gia tăng giá trị tại một trạm điển hình hay không. Cả hai đều được báo cáo xuyên suốt luận văn này.

\subsection{Đánh giá độc lập}

Một phép kiểm tra sâu hơn về khả năng chuyển giao của mô hình là đánh giá độc lập trên một bộ dữ liệu hoàn toàn không được dùng trong bất kỳ khía cạnh nào của quá trình phát triển mô hình. Trong luận văn này, đánh giá độc lập được thực hiện bằng 39 trạm cảm biến chi phí thấp (LCS) và một trạm tham chiếu Đại sứ quán Hoa Kỳ. Các trạm này được triển khai sau khi mô hình chính được phát triển, sử dụng phần cứng cảm biến khác (lớp Plantower PMS cho LCS; Met One BAM-1022 cho Đại sứ quán), và không được dùng cho kỹ thuật đặc trưng, tinh chỉnh siêu tham số, hay bất kỳ quyết định lựa chọn mô hình nào. Tuy nhiên, cần lưu ý rằng đánh giá độc lập này không phải là phép kiểm tra dự đoán không gian thực sự tại các vị trí hoàn toàn chưa thấy: các đặc trưng vệ tinh, AOD, TROPOMI, và khí tượng ERA5 được lấy từ bản ghi theo giờ của trạm quan trắc (KK) gần nhất, trong khi chỉ có PM2.5 mục tiêu đến từ cảm biến chi phí thấp giữ lại, và các điểm neo không gian RFSI được định vị đúng tại từng vị trí độc lập. Do đó tốt nhất nên hiểu nó là một phép kiểm tra chuyển đặc trưng từ trạm gần nhất cộng với sự phù hợp cảm biến, và sự suy giảm hiệu suất theo khoảng cách được báo cáo về sau phần nào phản ánh sai số thay thế đặc trưng tăng dần theo khoảng cách.

Đánh giá độc lập tránh được không chỉ các dạng rò rỉ có thể xảy ra trong LOSO --- ví dụ, nếu dữ liệu của trạm giữ lại đã ảnh hưởng đến các quyết định lựa chọn đặc trưng, lựa chọn siêu tham số, hay quyết định đưa vào một số nguồn dữ liệu nhất định --- mà còn cung cấp một phép đánh giá trên phần cứng cảm biến khác và giai đoạn thời gian khác. Cái giá phải đánh đổi là đánh giá độc lập chỉ khả thi khi tồn tại một mạng lưới quan trắc độc lập, vốn là một sự xa xỉ thực tiễn mà ít nghiên cứu có thể tận dụng.


Chương này đã trình bày các nền tảng lý thuyết cho luận văn: sự phân biệt giữa kỳ vọng ban đầu (nội suy thời gian) và bài toán trọng tâm (ngoại suy không gian) cùng các phương pháp đánh giá tương ứng, hiện trạng nghiên cứu trong ước tính PM2.5 dựa trên vệ tinh, các thách thức vật lý của mối quan hệ AOD-PM2.5 trong môi trường nhiệt đới, và khung mô hình hóa XGBoost/DART. Chương tiếp theo trình bày các nguồn dữ liệu cụ thể, quy trình kỹ thuật đặc trưng, và phương pháp đề xuất.
